\documentclass[a4paper]{article}
\usepackage[top=7mm,bottom=7mm,left=7mm,right=7mm]{geometry}
\usepackage[T1]{fontenc}
\usepackage{lmodern,amsmath,amssymb,mathtools,microtype,xcolor}
\usepackage{enumitem}
\renewcommand{\familydefault}{\sfdefault}
\pagestyle{empty}
\setlength{\parindent}{0pt}
\setlength{\parskip}{1.0pt}
\setlength{\jot}{1.1pt}
\definecolor{ink}{HTML}{152E3A}
\definecolor{muted}{HTML}{51616A}
\definecolor{line}{HTML}{B7C1C6}
\newcommand{\E}{\mathbb E}
\newcommand{\R}{\mathbb R}
\newcommand{\1}{\mathbf1}
\newcommand{\norm}[1]{\lVert#1\rVert}
\newcommand{\ip}[2]{\langle#1,#2\rangle}
\newcommand{\Var}{\operatorname{Var}}
\newcommand{\argmin}{\operatorname*{arg\,min}}
\newcommand{\argmax}{\operatorname*{arg\,max}}
\newcommand{\diag}{\operatorname{diag}}
\newcommand{\sect}[1]{\par\vspace{2.8pt}{\color{ink}\fontsize{8.6}{9.7}\selectfont\bfseries #1}\par\vspace{1.1pt}}
\newcommand{\kw}[1]{\textbf{#1}}
\newcommand{\smallnote}[1]{{\color{muted}#1}}
\newlength{\cw}\setlength{\cw}{\dimexpr(\textwidth-8mm)/3\relax}
\newsavebox{\A}\newsavebox{\B}\newsavebox{\C}\newsavebox{\D}\newsavebox{\F}\newsavebox{\G}
\newenvironment{sheetcol}[1]{\def\currentbox{#1}\begin{lrbox}{#1}\begin{minipage}[t]{\cw}\vspace{0pt}\fontsize{7.0}{8.0}\selectfont\raggedright\setlength{\abovedisplayskip}{1.5pt}\setlength{\belowdisplayskip}{1.5pt}\setlength{\abovedisplayshortskip}{1pt}\setlength{\belowdisplayshortskip}{1pt}}{\end{minipage}\end{lrbox}\global\setbox\currentbox=\box\currentbox}
\newcommand{\head}[2]{\noindent{\color{ink}\fontsize{15}{16}\selectfont\bfseries COMP3670}\hspace{6pt}{\fontsize{9}{11}\selectfont MACHINE LEARNING}\hfill{\color{muted}\fontsize{8}{10}\selectfont #1 / 2}\par\vspace{1.5pt}{\color{muted}\fontsize{7.1}{8.5}\selectfont #2}\par\vspace{3pt}{\color{ink}\hrule height .7pt}\vspace{3pt}}
\newcommand{\foot}[1]{\vfill\vspace{3pt}{\color{line}\hrule height .35pt}\vspace{2pt}{\color{muted}\fontsize{6.3}{7.2}\selectfont #1\hfill A4 / 100\% / duplex, long-edge}}
\begin{document}

\begin{sheetcol}{\A}
\sect{Learning from data}
\kw{Deduction:} theory/premises $\to$ prediction; true premises and valid logic imply the conclusion. \kw{Induction:} observations $\to$ a general rule; finite evidence cannot prove a universal claim. One counterexample can falsify it.

\kw{Supervised:} labelled $(x_i,y_i)$; classification (discrete $y$), regression (numerical $y$). \kw{Unsupervised:} only $x_i$; discover structure. \kw{Reinforcement:} state $\to$ action $\to$ reward; learn a policy maximizing cumulative reward, often delayed.

\kw{PLA} ($y_i\in\{-1,+1\}$). Initialize $w=b=0$. Find a misclassified point, update, repeat until no errors:
\[w\gets w+y_ix_i,\qquad b\gets b+y_i.\]
L1 predicts $+1$ if $w^Tx+b>0$, otherwise $-1$. Strictly linearly separable data $\Rightarrow$ finitely many updates; otherwise no termination guarantee. Zero training error does not imply generalization.

\sect{Variance \& numerical integration}
\[\Var(X)=\E[(X-\E X)^2]=\E[X^2]-(\E X)^2.\]
For $I=\int_a^b f(x)dx$, let $h=(b-a)/n$, $x_i=a+ih$, $f_i=f(x_i)$.

\kw{Composite trapezoidal rule}
\[T_n=\frac h2\left[f_0+2\sum_{i=1}^{n-1}f_i+f_n\right].\]
If $\max_{[a,b]}|f''|\le M_2$:
\[|I-T_n|\le\frac{M_2(b-a)^3}{12n^2}.\]
\kw{Composite Simpson rule} ($n$ even)
\[S_n=\frac h3\left[f_0+4\!\sum_{\substack{1\le i<n\\i\text{ odd}}}\!f_i
+2\!\sum_{\substack{2\le i<n\\i\text{ even}}}\!f_i+f_n\right].\]
If $\max|f^{(4)}|\le M_4$, $|I-S_n|\le M_4(b-a)^5/(180n^4)$. Doubling $n$: error bounds divide by $4$ (trapezoid), $16$ (Simpson).

\sect{Generalization \& complexity}
$D=\{(X_i,Y_i)\}_{i=1}^N$ i.i.d.; $L(h)=E_{out}(h)$, $\widehat L(h)=E_{in}(h)$; ERM minimizes $\widehat L$. Bounds below use binary 0--1 loss.

\kw{Hoeffding}, fixed $h$ before viewing $D$:
\[\Pr\{|L(h)-\widehat L(h)|>\epsilon\}\le2e^{-2N\epsilon^2}.\]
For finite $|\mathcal H|=M$, union bound gives
\[\Pr\{\sup_h|L(h)-\widehat L(h)|>\epsilon\}\le2Me^{-2N\epsilon^2}.\]
Thus simultaneously $L(h)\le\widehat L(h)+\sqrt{\log(2M/\delta)/(2N)}$, with probability $\ge1-\delta$. Searching over data-selected $h$ requires uniform control.

\kw{Shatter:} realize all $2^N$ labellings of some $N$ points. $S(\mathcal H,N)$: maximum number of labellings; $v=d_{VC}$: largest shatterable size. One-direction threshold: $1$; intervals: $2$; affine hyperplanes in $\R^d$: $d+1$.

\kw{Sauer:} finite $v$, $N>v$ $\Rightarrow S(\mathcal H,N)\le(N+1)^v$. Course VC bound:
\[\Pr\{\sup_h|L-\widehat L|>\epsilon\}\le8(N+1)^ve^{-N\epsilon^2/32}.\]
Penalty $\sqrt{\frac{32}{N}[v\log(N+1)+\log(8/\delta)]}$.

\kw{Agnostic PAC:} $\Pr\{L(\hat h)\le L(h^*)+\epsilon\}\ge1-\delta$ for all $P$, once $N\ge N_{\mathcal H}(\epsilon,\delta)$; $h^*=\argmin_{h\in\mathcal H}L(h)$. Binary distribution-free PAC $\Leftrightarrow v<\infty$ (standard measurability assumptions). Sufficient course-bound order:
\[N_{\mathcal H}=O\!\left(\frac{v\log(1/\epsilon)+\log(1/\delta)}{\epsilon^2}\right).\]
\kw{Approximation vs estimation:}
\[\underbrace{L(\hat h)-L(f^*)}_{\text{excess over Bayes}}=
\underbrace{L(h^*)-L(f^*)}_{\Delta_{app}}+
\underbrace{L(\hat h)-L(h^*)}_{\Delta_{est}}.\]
Restricting $\mathcal H$ can increase approximation error but reduce estimation difficulty. This is related to the bias/variance trade-off.

\kw{$K$-fold CV:} fit on $K-1$ folds, score the held-out fold; rotate, then average:
\[E_{CV}(M)=K^{-1}\sum_{j=1}^K E_{val}^{(j)}(M).\]
Choose hyperparameters by CV, refit on development data; test once on untouched data. Fit preprocessing inside each fold. Fold scores are dependent; CV winner's score can be optimistic.
\end{sheetcol}

\begin{sheetcol}{\B}
\sect{Parametric estimation \& logistic models}
\kw{Parametric:} fixed finite-dimensional family. \kw{Nonparametric:} complexity may grow with data; locality/smoothing replaces a fixed shape assumption. Plug-in: estimate $\eta$, then apply a decision rule.

\kw{Likelihood / MLE} (i.i.d. observations):
\[\mathcal L(\theta)=\prod_{i=1}^N p(z_i\mid\theta),\]
\[\hat\theta=\argmax_\theta\sum_i\log p(z_i\mid\theta).\]
Data fixed, parameter varies; likelihood is not a probability distribution over $\theta$. Bernoulli: $\hat p=k/N$. Gaussian MLE:
\[\hat\mu=N^{-1}\sum_i z_i,\qquad\hat\sigma^2=N^{-1}\sum_i(z_i-\hat\mu)^2.\]
Consistency/efficiency need correct specification and regularity; a wrong family can fit well yet remain wrong.

\kw{Logistic regression}, $y\in\{0,1\}$:
\[g=w^Tx+b,\quad p=\sigma(g)=\frac1{1+e^{-g}},\]
\[\log\frac p{1-p}=g.\]
Odds $=e^g$; increasing $x_j$ by one multiplies odds by $e^{w_j}$, holding others fixed. Threshold $p>1/2$ iff $g>0$. Costs can change the threshold.
\[\ell_{CE}(y,p)=-y\log p-(1-y)\log(1-p).\]
Conditional likelihood $\prod_i p_i^{y_i}(1-p_i)^{1-y_i}$; MLE $=$ minimize mean cross-entropy, not directly 0--1 error. Confident wrong probabilities incur large loss.

\kw{Logistic gradient} (augmented $x$ includes bias):
\[\sigma'(z)=\sigma(z)(1-\sigma(z)),\]
\[\nabla_w\ell_i=(\sigma(w^Tx_i)-y_i)x_i,\]
\[\nabla J=\frac1N\sum_i(p_i-y_i)x_i,\qquad
\nabla^2J=\frac1NX^TDX,\]
$D=\diag(p_i(1-p_i))\succeq0$. Convex, usually no closed form; unregularized separable data may have no finite MLE.

\sect{Nonparametric classification}
\kw{Histogram density}, bin $A(x)$, width $l$ in 1D:
\[\hat f(x)=\frac{\#\{i:X_i\in A(x)\}}{Nl}
.\]
In $d$ dimensions, divide by bin volume instead of $l$.
\kw{Kernel density estimate}, normalized $\kappa$:
\[\hat f(x)=\frac1{N\sigma^d}\sum_i\kappa\!\left(\frac{x-X_i}{\sigma}\right).\]
Small bins/bandwidth: high variance; large: more smoothing/bias. Scaling and irrelevant features affect distances; high dimension causes sparse neighborhoods.

\kw{Unified local vote} ($W_i\ge0$, $\sum_iW_i=1$):
\[\hat\eta(x)=\sum_i W_i(x)\1\{Y_i=1\},\]
\[\hat h(x)=\1\{\hat\eta(x)>1/2\}.\]
\kw{Histogram:} $W_i=\1\{X_i\in A(x)\}/N(x)$; handle empty cells. \kw{kNN:} $W_i=1/k$ for the $k$ nearest points, else $0$. Small $k$: flexible/noisy. Odd $k$ prevents binary vote ties, not distance ties.
\[\text{Kernel:}\quad W_i(x)=
\frac{\kappa((x-X_i)/\sigma)}{\sum_j\kappa((x-X_j)/\sigma)}.\]
Handle a zero denominator for compact-support kernels. Gaussian $e^{-\norm{u}^2}$; Cauchy $(1+\norm{u}^{d+1})^{-1}$; uniform $\1_{\norm{u}\le1}$; triangular $(1-\norm{u})_+$; Epanechnikov $(1-\norm{u}^2)_+$. Common constants cancel in normalized weights. Select $k,l,\sigma$ by validation/CV.

\sect{Trees, bagging \& random forests}
\kw{CART:} recursively split $x_j\le a$; leaves form axis-aligned boxes. Predict majority class in each leaf. Each path is an AND; same-class leaves are OR'ed.

Let $p$ be class-1 fraction in node $R$:
\[I_G=2p(1-p),\quad I_{err}=\min(p,1-p),\]
\[I_H=-p\log p-(1-p)\log(1-p),\quad0\log0=0.\]
Impurity $0$ at pure nodes, maximal at $p=1/2$.
\[\Delta I=I(R)-\frac{N_L}{N_R}I(R_L)-\frac{N_U}{N_R}I(R_U).\]
Here $N_R$ is parent size, $N_L,N_U$ are child sizes. Choose the largest admissible decrease; repeat. A local greedy split is not global ERM over all trees. Stop by purity, depth, minimum leaf size; prune/validate. Deep single trees have high variance. Monotone feature rescaling preserves threshold order.

\end{sheetcol}

\begin{sheetcol}{\C}
\sect{K-means}
For nonempty clusters $C_1,\ldots,C_k$:
\[J=\sum_{r=1}^k\sum_{i\in C_r}\norm{x_i-c_r}^2,\qquad
c_r=\frac1{|C_r|}\sum_{i\in C_r}x_i.\]
\kw{Lloyd steps:} initialize centers/partition; assign each point to its nearest center; recompute means; repeat until assignments stabilize. Use multiple starts; different initializations can reach different local optima.
\[\sum_i\norm{x_i-c}^2=
\sum_i\norm{x_i-\bar x}^2+n\norm{c-\bar x}^2.\]
Thus the mean uniquely minimizes a nonempty cluster's SSE. Assignment and mean updates never increase $J$. With no empty clusters and ties kept in the current cluster, any assignment change strictly decreases $J$; finitely many partitions $\Rightarrow$ finite termination, not global optimality.

\kw{Empty cluster:} its mean is undefined; repair/reseed (e.g. a distant point), or keep the previous center under a chosen policy. The stated termination proof assumes nonempty clusters, so do not silently divide by zero. Feature scaling changes the geometry.

\sect{Hierarchical clustering}
\kw{Agglomerative:} start with singleton clusters; repeatedly merge the closest pair under a chosen linkage; record merge height; continue to one cluster.
\[d_{single}(A,B)=\min_{x\in A,y\in B}d(x,y),\]
\[d_{complete}(A,B)=\max_{x\in A,y\in B}d(x,y),\]
\[d_{average}(A,B)=\frac1{|A||B|}\sum_{x\in A}\sum_{y\in B}d(x,y).\]
Single can chain via bridges; complete favors compact clusters; average uses all cross-pairs. A dendrogram's vertical height is merge dissimilarity. A horizontal cut gives clusters as connected branches below it; horizontal spacing has no distance meaning. Linkage choice can change the partition at the same cut height.

\sect{Orthonormal geometry}
\kw{Cauchy--Schwarz:}
\[|x^Ty|\le\norm{x}\norm{y}.\]
Equality iff $x,y$ linearly dependent (including a zero vector). Proof: $0\le\norm{x-ay}^2$, choose $a=x^Ty/\norm y^2$ for $y\ne0$.

\kw{Orthonormal list:} $q_i^Tq_j=\delta_{ij}$, hence linearly independent. If it spans $S$, it is an orthonormal basis. With $Q=[q_1\ \cdots\ q_k]$:
\[Q^TQ=I_k,\qquad [x]_{\mathcal Q}=Q^Tx\quad(x\in S).\]
The coefficients follow by taking dot products. Rectangular $Q$ generally has $QQ^T\ne I_d$.

\kw{Orthogonal matrix} is square with $Q^TQ=I$:
\[Q^{-1}=Q^T,\quad (Qx)^T(Qy)=x^Ty,\quad\norm{Qx}=\norm x.\]
It preserves lengths, distances and nonzero-vector angles. Do not confuse orthogonal columns with a square orthogonal matrix.

\kw{Gram--Schmidt}, independent $v_1,\ldots,v_k$:
\[u_1=v_1,\quad q_1=u_1/\norm{u_1},\]
\[u_j=v_j-\sum_{i<j}(q_i^Tv_j)q_i,\qquad
q_j=u_j/\norm{u_j}.\]
Subtract \emph{all} previous components, then normalize. Prefix spans are unchanged. Zero residual means dependence; skip it rather than normalize zero. Check $Q^TQ=I$ afterwards.

\sect{Spectral theorem \& PSD}
For real symmetric $A=A^T\in\R^{d\times d}$, there is a real orthonormal eigenbasis:
\[Aq_j=\lambda_jq_j,\quad Q^TQ=I,\]
\[A=Q\Lambda Q^T=\sum_{j=1}^d\lambda_jq_jq_j^T.\]
Distinct-eigenvalue eigenvectors are orthogonal; within a repeated eigenspace choose an orthonormal basis. $Q^T$: obtain coordinates; $\Lambda$: scale; $Q$: reconstruct.
\[x=\sum_j a_jq_j\quad\Rightarrow\quad x^TAx=\sum_j\lambda_ja_j^2.\]
\kw{Positive semidefinite:} $A\succeq0$ iff $x^TAx\ge0$ for all $x$, iff every $\lambda_j\ge0$. \kw{Positive definite:} strict positivity for every nonzero $x$, iff every $\lambda_j>0$.

PSD can be singular; PD is invertible. For any real $B$, $B^TB\succeq0$ since $x^TB^TBx=\norm{Bx}^2$. Negative eigenvalues rule out PSD. A general real non-symmetric matrix need not have an orthonormal eigenbasis.
\end{sheetcol}

\begin{sheetcol}{\D}
\sect{SVM: hard, soft \& nonlinear}
Labels $y_i\in\{-1,+1\}$; $f_i=w^Tx_i+b$. SVM convention: predict $+1$ for $f\ge0$, else $-1$. With $w\ne0$:
\[m_i=y_if_i,\quad\gamma_i=\frac{m_i}{\norm w},\quad
d(x,H)=\frac{|w^Tx+b|}{\norm w}.\]
Positive rescaling preserves boundary/margins, not raw scores.

\kw{Hard margin} (both classes, strictly separable):
\[\min_{w,b}\tfrac12\norm w^2\quad\text{s.t. }m_i\ge1\ (\forall i).\]
At optimum $\min_i m_i=1$, training margin $\gamma=1/\norm w$, full band width $2/\norm w$. \kw{Support vectors:} points with $m_i=1$ touch the margin boundaries. 

\kw{Soft margin}, $C>0$:
\[\begin{aligned}
\min_{w,b,\xi}\quad&\tfrac12\norm w^2+C\sum_i\xi_i,\\
\text{s.t.}\quad&m_i\ge1-\xi_i,\quad\xi_i\ge0.
\end{aligned}\]
For fixed $w,b$, $\xi_i^*=(1-m_i)_+=\max(0,1-m_i)$; equivalent objective:
\[J(w,b)=\tfrac12\norm w^2+C\sum_i(1-m_i)_+.\]
$\xi=0$: satisfies margin; $0<\xi<1$: correct side inside band; $\xi=1$: decision boundary; $\xi>1$: wrong side. Correct labels may still incur hinge loss. In soft SVM, inside-margin/violating points can influence the fit too.

Large $C$: penalize violations more; small $C$: prefer small norms. Feature units change distances and regularization geometry; tune scaling and $C$ together. Soft margin allows $w=0$, a constant score.

\kw{QP:} quadratic convex objective plus linear constraints; every local minimum is global. 

\kw{Subgradient step.} Compute $I=\{i:m_i<1\}$ from the same old parameters, then update both:
\[\begin{aligned}
w^+&=(1-\eta)w+\eta C\sum_{i\in I}y_ix_i,\\
b^+&=b+\eta C\sum_{i\in I}y_i.
\end{aligned}\]
Shrink weights and correct violations. At $m=1$, zero hinge subgradient is valid. A step may increase $J$; arbitrary constant steps need not converge.

\kw{Nonlinear SVM:} choose a fixed feature map $\phi$; replace $x_i$ by $\phi(x_i)$ in the score, hinge objective and updates. Still convex in $(w,b)$; curved boundary in original space. E.g. $\phi(x)=(x,x^2)$, $4x-x^2-3>0\iff1<x<3$. A radial feature $x_1^2+x_2^2$ represents circular boundaries.

\sect{SVD \& latent-factor recommendation}
$A\in\R^{m\times d}$, rank $r$; compact SVD:
\[A=U_r\Sigma_rV_r^T=\sum_{j=1}^r\sigma_ju_jv_j^T,\]
$U_r\in\R^{m\times r},V_r\in\R^{d\times r}$ have orthonormal columns; $\sigma_1\ge\cdots\ge\sigma_r>0$.
\[Av_j=\sigma_ju_j,\quad A^Tu_j=\sigma_jv_j,\quad A^TAv_j=\sigma_j^2v_j.\]
\kw{Truncation:} $A_k=U_k\Sigma_kV_k^T$ is best among rank-$\le k$ matrices in Frobenius error:
\[\norm{A-A_k}_F^2=\sum_{j>k}\sigma_j^2,\quad
\norm A_F^2=\sum_j\sigma_j^2.\]
\kw{Ratings:} $R\in\R^{n\times m}$, $P\in\R^{n\times k},Q\in\R^{m\times k}$; $\hat r_{ij}=p_i^Tq_j$, $\hat R=PQ^T$. Fully observed: $P=U_k\Sigma_k,Q=V_k$. Rows are individual embeddings; columns are patterns.

\kw{Missing ratings:} unknown $\ne0$. Fit only observed training pairs $\Omega$:
\[L=\sum_{(i,j)\in\Omega}(r_{ij}-p_i^Tq_j)^2+\lambda(\norm P_F^2+\norm Q_F^2).\]
\kw{ALS} ($\lambda>0$): small random initialization; alternate exact ridge solves. Let $Q_i$ contain item-factor rows rated by user $i$, $r_i$ their ratings; $P_j$ contains users rating item $j$, with ratings $r_j$:
\[p_i=(Q_i^TQ_i+\lambda I)^{-1}Q_i^Tr_i,\]
\[q_j=(P_j^TP_j+\lambda I)^{-1}P_j^Tr_j.\]
Solve systems, do not explicitly invert. Exact blocks do not increase $L$; joint problem nonconvex, no global guarantee. Validate $k,\lambda$; untouched test ratings assess prediction. No training records: cold start.
\[\mathrm{RMSE}=\sqrt{\frac1{|T|}\sum_{(i,j)\in T}(r_{ij}-\hat r_{ij})^2}.\]
\end{sheetcol}

\begin{sheetcol}{\F}
\sect{Second-order tests \& convexity}
\kw{1D stationary point:} $f'(x^*)=0$ (interior, differentiable) is necessary, not sufficient. If $f''(x^*)>0$: strict local min; $<0$: strict local max; $=0$: inconclusive. Check boundaries and non-differentiable points too.

\kw{Multivariable:} $\nabla f(x^*)=0$; Hessian $H_{ij}=\partial_i\partial_j f$ (symmetric for $C^2$ functions).
\[f(x+v)\approx f(x)+\nabla f(x)^Tv+\tfrac12v^TH(x)v.\]
At a stationary point: $H\succ0$ implies strict local min; $H\prec0$ implies strict local max; indefinite $H$ implies saddle. Semidefinite/singular $H$: test inconclusive. Local min requires $H\succeq0$ for $C^2$ functions.

\kw{2D test} at a stationary point, $D=f_{xx}f_{yy}-f_{xy}^2$: $D>0,f_{xx}>0$ gives a min; $D>0,f_{xx}<0$ a max; $D<0$ a saddle; $D=0$ is inconclusive.

\kw{Convexity} on a convex domain:
\[f(tx+(1-t)y)\le tf(x)+(1-t)f(y),\quad t\in[0,1].\]
For differentiable $f$, equivalent to $f(y)\ge f(x)+\nabla f(x)^T(y-x)$. For $C^2$ $f$, equivalent to $H(x)\succeq0$ everywhere. Every local min is global; $\nabla f=0$ is sufficient for global min. Strict convexity gives at most one minimizer, not its existence. $H\succ0$ everywhere is sufficient for strict convexity, not necessary.

\sect{Normal equations \& ridge}
Rows of design $X\in\R^{N\times p}$ include any intercept/basis features; $y\in\R^N$.
\[RSS=\norm{X\theta-y}^2,\quad
\nabla RSS=2X^T(X\theta-y),\]
\[\nabla^2RSS=2X^TX\succeq0,\qquad X^TX\hat\theta=X^Ty.\]
If $X$ has full column rank:
\[\hat\theta=(X^TX)^{-1}X^Ty.\]
Otherwise minimizers need not be unique; fitted training predictions are unique. Numerically solve least squares, not an explicit inverse.

\kw{Ridge} with $\lambda>0$:
\[J=RSS+\lambda\norm\theta^2,\quad
\hat\theta=(X^TX+\lambda I)^{-1}X^Ty,\]
\[\nabla J=2X^T(X\theta-y)+2\lambda\theta.\]
$X^TX+\lambda I\succ0$: unique solution and improved conditioning. Formula penalizes all entries; to leave intercept unpenalized replace $I$ by $\diag(0,1,\ldots,1)$ (uniqueness then needs corresponding conditions).

\kw{Parametric vs local regression.} Linear parameter models use $h_\theta(x)=\sum_j\theta_j\phi_j(x)$; features may be nonlinear in $x$. For squared loss, $f^*(x)=\E[Y\mid X=x]$.
\kw{Kernel / kNN regression:} $\hat h(x)=\sum_iW_i(x)Y_i$, using normalized kernel weights or $1/k$ for nearest neighbors. \kw{Regression trees:} leaf mean, splits minimizing child RSS; forests average predictions.

\sect{GD, learning rates \& convergence}
\[\theta_{t+1}=\theta_t-\alpha\nabla J(\theta_t).\]
Full GD uses all $N$ examples per step. For $J=RSS/N$:
\[\theta^+=\theta-\alpha\frac2N X^T(X\theta-y).\]
Large $\alpha$ may oscillate/diverge; tiny $\alpha$ is slow. Stop by small gradient, small objective/parameter changes, or iteration budget; validation can guide early stopping.

\kw{Convex $K$-smooth objective}, minimizer $\theta^*$ exists and $0<\alpha\le1/K$:
\[J(\theta_t)-J(\theta^*)\le\frac{\norm{\theta_0-\theta^*}^2}{2\alpha t}.\]
$K$ bounds Hessian eigenvalues / gradient Lipschitz constant. For strongly convex $J$ with curvature $\mu>0$, suitable GD has a linear rate. Nonconvex descent need not find a global minimum.

\kw{SGD / mini-batches}, $J=N^{-1}\sum_i\ell_i$:
\[g_t=\frac1b\sum_{i\in B_t}\nabla\ell_i(\theta_t),\quad
\theta_{t+1}=\theta_t-\alpha_tg_t.\]
Uniform sampling gives $\E[g_t\mid\theta_t]=\nabla J(\theta_t)$. One epoch $\approx N/b$ updates. Smaller $b$: cheaper, noisier steps; larger $b$: less gradient variance, more cost. Apply regularizer gradients separately, without dropping their scaling.

A fixed learning rate often leaves a noise floor. Under standard stochastic-approximation assumptions, a sufficient diminishing-step condition is
\[\sum_t\alpha_t=\infty,\qquad\sum_t\alpha_t^2<\infty.\]
E.g. $\alpha_t=a/(b+t)^\gamma$, $1/2<\gamma\le1$. These step rules alone do not guarantee convergence without assumptions on the objective and gradient noise.
\end{sheetcol}

\begin{sheetcol}{\G}
\sect{Ensemble learning}
\kw{Bagging:} draw $B$ bootstrap samples of size $N$ with replacement; fit a tree per sample; majority vote (classification), average (regression). Variance reduction is greater when members are less correlated.

\kw{Random forest:} bootstrap plus a random feature subset \emph{at every split}; often deep trees. Feature randomness decorrelates trees. Forests are more stable, less interpretable; no universal immunity to overfitting.

\sect{Neural networks \& backpropagation}
\kw{Forward pass:} $a^0=x$; hidden layers
\[z^\ell=W^\ell a^{\ell-1}+b^\ell,\qquad
 a^\ell=\phi(z^\ell).\]
Without nonlinear activations, stacked affine layers collapse to one affine map. Training is generally nonconvex; expressiveness does not guarantee generalization.

\kw{Activations:} ReLU $(z)_+=\max(0,z)$; $\sigma(z)=1/(1+e^{-z})$; $\tanh z$; $\arctan z$; step $\1\{z\ge0\}$.
\[\phi'_{ReLU}=\1\{z>0\}\ (z\ne0),\quad
\sigma'=\sigma(1-\sigma),\]
\[(\tanh z)'=1-\tanh^2z,\quad (\arctan z)'=(1+z^2)^{-1}.\]
ReLU at zero: choose a subgradient (often 0); step has no useful ordinary gradient at its jump.

\kw{Output:} binary sigmoid + binary cross-entropy; multiclass softmax + cross-entropy:
\[p_c=\frac{e^{z_c}}{\sum_j e^{z_j}},\quad \hat y=\argmax_c p_c,\]
\[\ell=-\sum_c y_c\log p_c.\]
Use one-hot $y$; subtract $\max_jz_j$ before exponentiating for numerical stability. For either matched sigmoid/BCE or softmax/CE, output error is $\delta^L=p-y$.

\kw{Backprop} (linear output logits):
\[\delta^\ell=((W^{\ell+1})^T\delta^{\ell+1})\odot\phi'(z^\ell),\]
\[\frac{\partial\ell}{\partial W^\ell}=\delta^\ell(a^{\ell-1})^T,\qquad
\frac{\partial\ell}{\partial b^\ell}=\delta^\ell.\]
Average sample gradients for mean batch loss. Chain rule transmits error backwards; keep tensor shapes consistent. Saturating activations can yield small gradients.

\kw{Initialization:} small random zero-mean weights break neuron symmetry; biases may be zero. Weight variances (not standard deviations), fan-in $n_{in}$, fan-out $n_{out}$:
\[\text{LeCun }\frac1{n_{in}},\quad
\text{Xavier }\frac2{n_{in}+n_{out}},\quad
\text{He }\frac2{n_{in}}.\]
He suits ReLU; Xavier balances layer scales. All-zero hidden weights keep units symmetric.

\sect{Optimizer update equations}
Let $g_t=\nabla J(w_t)$ or a mini-batch gradient; all vector squares, roots, products/divisions below are coordinatewise. Initialize historical states to zero. $0<\beta,\beta_1,\beta_2<1$; $\epsilon>0$ stabilizes denominators.

\kw{Momentum / heavy ball} (velocity convention):
\[v_t=\beta v_{t-1}-\alpha_tg_t,\qquad w_{t+1}=w_t+v_t.\]
Accumulates descent directions; can reduce oscillation.

\kw{Nesterov} (gradient at look-ahead position):
\[v_t=\beta v_{t-1}-\alpha_t\nabla J(w_t+\beta v_{t-1}),\]
\[w_{t+1}=w_t+v_t.\]
Evaluate the gradient at the look-ahead position.

\kw{Adagrad} (permanent squared-gradient sum):
\[s_t=s_{t-1}+g_t\odot g_t,\qquad
w_{t+1}=w_t-\frac{\alpha_tg_t}{\sqrt{s_t+\epsilon}}.\]
Effective coordinatewise learning rates shrink; useful for sparse coordinates, but may become too small.

\kw{RMSProp} (EMA instead of permanent sum):
\[s_t=\beta s_{t-1}+(1-\beta)g_t\odot g_t,\]
\[w_{t+1}=w_t-\frac{\alpha_tg_t}{\sqrt{s_t+\epsilon}}.\]
Forgets old squared gradients. Course convention puts $\epsilon$ inside the root in these two methods.

\kw{Adam} (course's \emph{negative-gradient} moment):
\[m_t=\beta_1m_{t-1}-(1-\beta_1)g_t,\]
\[s_t=\beta_2s_{t-1}+(1-\beta_2)g_t\odot g_t,\]
\[\hat m_t=\frac{m_t}{1-\beta_1^t},\qquad
\hat s_t=\frac{s_t}{1-\beta_2^t},\qquad t\ge1,\]
\[w_{t+1}=w_t+\frac{\alpha_t\hat m_t}{\sqrt{\hat s_t+\epsilon_{root}}+\epsilon}.\]
Use $m_0=s_0=0$; first gradient is $t=1$ for bias correction. Common course values: $\beta_1=.9$, $\beta_2=.999$, $\epsilon=10^{-8}$, $\epsilon_{root}=0$.

\kw{Sign check:} negative gradient is already stored in $m_t$, so Adam's final update is \emph{plus}. With a positive-gradient first moment, use a minus update instead. Do not mix conventions. Bias correction compensates zero-initialized EMAs, not statistical model bias.
\end{sheetcol}

\typeout{COL-A=\the\dimexpr\ht\A+\dp\A\relax}
\typeout{COL-B=\the\dimexpr\ht\B+\dp\B\relax}
\typeout{COL-C=\the\dimexpr\ht\C+\dp\C\relax}
\typeout{COL-D=\the\dimexpr\ht\D+\dp\D\relax}
\typeout{COL-F=\the\dimexpr\ht\F+\dp\F\relax}
\typeout{COL-G=\the\dimexpr\ht\G+\dp\G\relax}
\head{1}{LEARNING / STATISTICS / MODELS / CLUSTERING / ORTHONORMAL GEOMETRY}
\noindent\usebox{\A}\hfill\usebox{\B}\hfill\usebox{\C}
\foot{Selected course notes / L1--4, 5b, 8a--b, 9--12 and LA5--8 / bounds assume their stated conditions}
\newpage
\head{2}{SVM / SVD / SECOND-ORDER TESTS / OPTIMISATION / NEURAL NETWORKS / OPTIMIZERS}
\noindent\usebox{\D}\hfill\usebox{\F}\hfill\usebox{\G}
\foot{Notation: unqualified norms are Euclidean / $I$ is identity / optimizer conventions follow Lecture 12}
\end{document}
